\begin{apartat}{1,25}
Escriviu el nombre màssic, el nombre atòmic i el nombre de neutrons que té el
$^{13}_{7}\mathrm{N}$. En la desintegració del $^{13}_{7}\mathrm{N}$, un dels protons del
nucli de nitrogen es transforma segons una desintegració
$\beta^+: {}^{1}_{1}\mathrm{p}^+ \to {}^{1}_{0}\mathrm{n} + {}^{0}_{+1}\mathrm{e}^+ + {}^{0}_{0}\nu$.
Escriviu la reacció de desintegració del $^{13}_{7}\mathrm{N}$. Justifiqueu per què la
desintegració $\beta^+$ no pot tenir lloc fora d'un nucli.
\end{apartat}

\begin{apartat}{1,25}
A partir de l'equació de l'evolució de la desintegració, determineu la relació entre la
constant de desintegració $\lambda$ i el període de semidesintegració. El període de
semidesintegració del $^{13}_{7}\mathrm{N}$ és de 9,965 min. Si en un instant determinat
hi ha una massa de 5 ng de $^{13}_{7}\mathrm{N}$, quina quantitat en romandrà al cap de
30 min?

Doneu l'expressió de l'activitat radioactiva en funció del temps. Quant de temps cal
esperar perquè l'activitat radioactiva es redueixi fins a un 1\,\% del seu valor inicial?
\end{apartat}

\dades{$_{4}\mathrm{Be}$\ $_{5}\mathrm{B}$\ $_{6}\mathrm{C}$\ $_{7}\mathrm{N}$\ $_{8}\mathrm{O}$\ $_{9}\mathrm{F}$\ $_{10}\mathrm{Ne}$\ $_{84}\mathrm{Po}$\ $_{85}\mathrm{At}$\ $_{86}\mathrm{Rn}$\ $_{87}\mathrm{Fr}$\ $_{88}\mathrm{Ra}$\ $_{89}\mathrm{Ac}$\ $_{90}\mathrm{Th}$.}
